% Szkielet sprawozdania LabSK — kopiuj do labsk/cX/sprawozdanie.tex
\documentclass[11pt,a4paper]{article}
\usepackage[utf8]{inputenc}
\usepackage[T1]{fontenc}
\usepackage[polish]{babel}
\usepackage{geometry}
\usepackage{graphicx}
\usepackage{listings}
\usepackage{xcolor}
\geometry{margin=2.2cm}
\lstset{
  basicstyle=\ttfamily\footnotesize,
  breaklines=true,
  frame=single,
  backgroundcolor=\color{gray!8},
  aboveskip=2pt,belowskip=2pt,
  xleftmargin=4pt,xrightmargin=4pt
}
% --- wypełnij ---
\newcommand{\autor}{Imię Nazwisko}
\newcommand{\album}{000000}
\newcommand{\cwiczenie}{X}
\newcommand{\temat}{Temat ćwiczenia}
\newcommand{\grupa}{gr. ...}

\title{LabSK\,/\,\cwiczenie: \temat}
\author{\autor \quad nr\,albumu:\,\album \quad \grupa}
\date{\today}

\begin{document}
\maketitle
\vspace{-1.5em}

% Schemat: \includegraphics[width=\linewidth]{ADM.svg}

\section*{Cel ćwiczenia}
% Skopiuj treść z Cel-ćwiczenia.txt — 2--3 zdania.

\section*{Zadania i dowody wykonania}
% Dla każdego zadania: 1 zdanie opisu + wywołanie + wynik (max.~10 linii).

\end{document}
